\chapter{正向极限和逆向极限}

记号约定:$I$是非空的预序集.

\section{原群的逆向系统}

\begin{definition}{原群的逆向极限}{}
	设$\left(\left(E_{\alpha}\right)_{\alpha\in I},\left(f_{\alpha\beta}\right)_{\alpha,\beta\in I}\right)$是集合的逆向系统.若
	\begin{enumerate}
		\item $\left(E_{\alpha}\right)_{\alpha\in I}$是一族原群,
		\item $\forall\alpha,\beta\in I\left(\alpha\leq\beta\implies f_{\alpha\beta}\in\Hom_{\mathsf{Mag}}\left(E_{\alpha},E_{\beta}\right)\right)$,
	\end{enumerate}
	则称$\left(\left(E_{\alpha}\right)_{\alpha\in I},\left(f_{\alpha\beta}\right)_{\alpha,\beta\in I}\right)$是原群的逆向系统,称$\varprojlim E_{\alpha}$是该逆向系统的逆向极限.
\end{definition}

\begin{remark}{}{}
	$\left(E_{\alpha}\right)_{\alpha\in I}$是$\prod\limits_{\alpha\in I}E_{\alpha}$的子原群.
\end{remark}

\begin{proposition}{原群的逆向极限的泛性质}{}
	设$\left(\left(E_{\alpha}\right)_{\alpha\in I},\left(f_{\alpha\beta}\right)_{\alpha,\beta\in I}\right)$是原群的逆向系统.
	\begin{enumerate}
		\item $\left(f_{\alpha}\right)_{\alpha\in I}$是一族原群同态.
		\item 设$F$是原群,对每个$\alpha\in I$有$u_{\alpha}\in\Hom_{\mathsf{Mag}}\left(F,E_{\alpha}\right)$.若$\forall\alpha,\beta\in I\left(\alpha\leq\beta\implies u_{\alpha}=f_{\alpha\beta}\circ u_{\beta}\right)$,则
		\begin{align*}
			\exists!u\in\Hom_{\mathsf{Mag}}\left(F,\varprojlim E_{\alpha}\right)\forall\alpha\in I\left(u_{\alpha}=f_{\alpha}\circ u\right).
		\end{align*}
	\end{enumerate}
\end{proposition}

\begin{remark}{}{}
	\begin{enumerate}
		\item 如果$\left(E_{\alpha}\right)_{\alpha\in I}$是一族结合原群(相对的,交换原群),则$\varprojlim E_{\alpha}$是结合原群(相对的,交换原群).
		\item 如果对每个$\alpha\in I$,原群$E_{\alpha}$的单位元是$e_{\alpha}$,并且$\forall\alpha,\beta\in I\left(\alpha\leq\beta\implies f_{\alpha\beta}\text{是幺原群同态}\right)$,则$\left(e_{\alpha}\right)_{\alpha\in I}$是$\varprojlim E_{\alpha}$的单位元,并且$\left(f_{\alpha}\right)_{\alpha\in I}$是一族幺原群同态.进一步,若$\left(u_{\alpha}\right)_{\alpha\in I}$是一族幺原群同态,则$u$也是要原群同态.
		\item 设$\left(x_{\alpha}\right)_{\alpha\in I}\in\varprojlim E_{\alpha}$,则$\left(x_{\alpha}\right)_{\alpha\in I}$可逆$\iff$对每个$\alpha\in I$,元素$x_{\alpha}$可逆.这时,$\left(x_{\alpha}\right)_{\alpha\in I}$的逆元为$\left(x_{\alpha}^{-1}\right)_{\alpha\in I}$.
		\item 如果$\left(E_{\alpha}\right)_{\alpha\in I}$是一族幺半群(相对的,群,环),$f_{\alpha\beta}(\alpha,\beta\in I,\alpha\leq\beta)$是幺半群(相对的,群,环)同态,那么$\varprojlim E_{\alpha}$是幺半群(相对的,群,环).类似的论断同样成立,在此我们略过相关的定义.
	\end{enumerate}
\end{remark}

\begin{proposition}{}{}
	设$\left(\left(E_{\alpha}\right)_{\alpha\in I},\left(f_{\alpha\beta}\right)_{\alpha,\beta\in I}\right),\left(\left(F_{\alpha}\right)_{\alpha\in I},\left(g_{\alpha\beta}\right)_{\alpha,\beta\in I}\right)$是原群逆系统,$\left(u_{\alpha}\right)_{\alpha\in I}$是前者到后者的同态,则$\varprojlim u_{\alpha}$是原群同态.
\end{proposition}

\begin{remark}{}{}
	对于幺半群,群,环我们有同样的论断,只需要把原群换成幺半群(群,环),原群同态换成幺半群(群,环)同态.
\end{remark}













